\documentclass[11pt]{article}
\usepackage[margin=1in]{geometry}
\usepackage{amsmath,amssymb,amsthm}
\usepackage{hyperref}

\newtheorem{theorem}{Theorem}
\newtheorem{lemma}{Lemma}
\newtheorem{remark}{Remark}

\newcommand{\R}{\mathbb R}
\newcommand{\Xclass}{\mathcal X}

\title{A decreasing-power Orlicz family outside the class \texorpdfstring{$\mathcal X$}{X}}
\author{Run fa\_banach\_001, agent lane 17}
\date{2026-07-03}

\begin{document}
\maketitle

\section*{Source and target}

The source paper is Astashkin--Lykov--Milman, \emph{Majorization revisited:
Comparison of norms in interpolation scales}, arXiv:2107.11854 \cite{ALM}.
In Section 6 the authors define the class $\Xclass$ of rearrangement invariant
spaces $X$ on $(0,\infty)$ with the following property. If a pair $(f,g)$
satisfies the Nazarov--Podkorytov single-crossing condition
\[
 f^*(t) \ge g^*(t)\quad (t\le t_0),\qquad
 f^*(t) \le g^*(t)\quad (t\ge t_0),
\]
and if $\|f\|_X\ge \|g\|_X$, then
\[
 \|f^*\chi_{(0,s)}\|_X\ge \|g^*\chi_{(0,s)}\|_X
 \qquad\text{for every }s>0.
\]
They state that a complete characterization of $\Xclass$ is open. They prove
that Lorentz spaces $\Lambda_r(\varphi)$ belong to $\Xclass$, that the
Marcinkiewicz spaces $M(t^\alpha)$ do not, and that an Orlicz space $L_M$
belongs to $\Xclass$ under a monotonicity assumption on
$M'(\varepsilon x)/M'(x)$.

\section*{Result}

This packet gives a negative Orlicz-family classification complementary to the
positive Orlicz theorem in the source paper.

\begin{theorem}
\label{thm:decreasing-power}
Let $0<p<q$. Define
\[
 N_{p,q}(u)=
 \begin{cases}
   u^q, & 0\le u\le 1,\\
   u^p, & u\ge 1,
 \end{cases}
 \qquad
 M_{p,q}(u)=\int_0^u N_{p,q}(v)\,dv .
\]
Then $M_{p,q}$ is a $\Delta_2$ Orlicz function, but the Orlicz space
$L_{M_{p,q}}(0,\infty)$ equipped with the Luxemburg norm is not in
$\Xclass$.
\end{theorem}

\section*{Proof intuition}

The source paper's positive Orlicz theorem says, in effect, that scaling an
Orlicz modular down by $\varepsilon$ must preserve the direction of a
single-crossing modular comparison. The family above has the opposite
geometry: near zero the modular behaves like a higher power, while at infinity
it behaves like a lower power. Thus scaling down suppresses the low-level
part more strongly than the high-level part.

We exploit this by building a two-step prefix. On the prefix, $f$ has a short
high block and a long low block, while $g$ has one intermediate block. At the
larger modular scale this prefix favors $g$; after scaling by $\varepsilon$ it
favors $f$. A common low tail is then appended to both functions. The common
tail does not affect the failed prefix comparison, but it normalizes the global
Luxemburg comparison so that $\|f\|\ge \|g\|$.

\section*{A crossing-block criterion}

We use the Luxemburg norm
\[
 \|h\|_{L_M}=\inf\left\{\lambda>0:
 \int_0^\infty M(|h(t)|/\lambda)\,dt\le 1\right\}.
\]

\begin{lemma}
\label{lem:block}
Let $M$ be a strictly increasing Orlicz function. Suppose there are
$0<d<a<b$ and $0<\varepsilon<1$ such that
\[
\frac{M(\varepsilon a)-M(\varepsilon d)}{M(a)-M(d)}
<
\frac{M(\varepsilon b)-M(\varepsilon a)}{M(b)-M(a)}.
\tag{1}
\label{eq:increment}
\]
Then $L_M(0,\infty)\notin \Xclass$.
\end{lemma}

\begin{proof}
Condition \eqref{eq:increment} is equivalent to the existence of a number
$\theta>0$ such that
\[
 \frac{M(b)-M(a)}{M(a)-M(d)}
 < \theta <
 \frac{M(\varepsilon b)-M(\varepsilon a)}
      {M(\varepsilon a)-M(\varepsilon d)} .
\tag{2}
\label{eq:theta}
\]
The left inequality gives
\[
 M(b)+\theta M(d)<(1+\theta)M(a).
\tag{3}
\label{eq:large-scale}
\]
Choose $B>0$ so that
\[
 B\bigl(M(b)+\theta M(d)\bigr)<1
 <B(1+\theta)M(a),
\tag{4}
\label{eq:B-choice}
\]
and put $C=\theta B$ and $s_0=B+C$.

On $(0,s_0)$ define decreasing step functions
\[
 U_0(t)=
 \begin{cases}
 b,&0<t<B,\\
 d,&B<t<s_0,
 \end{cases}
 \qquad
 V_0(t)=a\quad(0<t<s_0).
\]
By \eqref{eq:B-choice},
\[
 \int_0^{s_0}M(U_0(t))\,dt<1<
 \int_0^{s_0}M(V_0(t))\,dt.
\tag{5}
\label{eq:prefix-large}
\]
The right inequality in \eqref{eq:theta} gives
\[
 B M(\varepsilon b)+C M(\varepsilon d)
 >(B+C)M(\varepsilon a).
\tag{6}
\label{eq:small-scale}
\]
Since $\varepsilon<1$ and $M$ is increasing, the left side of
\eqref{eq:small-scale} is strictly smaller than
$B M(b)+C M(d)$, hence is $<1$ by \eqref{eq:B-choice}. Therefore
\[
 T=\frac{1-\int_0^{s_0}M(\varepsilon U_0(t))\,dt}{M(\varepsilon d)}
\]
is positive. Extend the unscaled functions by a common tail:
\[
 U(t)=
 \begin{cases}
 U_0(t),&0<t<s_0,\\
 d,&s_0<t<s_0+T,\\
 0,&t>s_0+T,
 \end{cases}
 \qquad
 V(t)=
 \begin{cases}
 V_0(t),&0<t<s_0,\\
 d,&s_0<t<s_0+T,\\
 0,&t>s_0+T.
 \end{cases}
\]
Finally set $f=\varepsilon U$ and $g=\varepsilon V$.

The functions are already decreasing rearrangements. With $t_0=B$ they satisfy
the source paper's $NP$ condition: $f^*\ge g^*$ on $(0,B)$ and
$f^*\le g^*$ on $(B,\infty)$.

At the global scale $\lambda=1$ we have, by the choice of $T$,
\[
 \int_0^\infty M(f(t))\,dt=1,
\]
while \eqref{eq:small-scale} implies
\[
 \int_0^\infty M(g(t))\,dt
 =
 \int_0^{s_0}M(\varepsilon V_0(t))\,dt+T M(\varepsilon d)
 <1.
\]
Thus $\|f\|_{L_M}=1\ge \|g\|_{L_M}$.

However, on the prefix $(0,s_0)$, \eqref{eq:prefix-large} says
\[
 \int_0^{s_0}M(f(t)/\varepsilon)\,dt<1
 <
 \int_0^{s_0}M(g(t)/\varepsilon)\,dt.
\]
Consequently
\[
 \|f^*\chi_{(0,s_0)}\|_{L_M}\le \varepsilon
 < \|g^*\chi_{(0,s_0)}\|_{L_M}.
\]
This contradicts the defining implication for $\Xclass$, so
$L_M\notin \Xclass$.
\end{proof}

\section*{Proof of the theorem}

The function $N_{p,q}$ is increasing and continuous, so $M_{p,q}$ is an Orlicz
function. Its growth is polynomial at both zero and infinity, so
$M_{p,q}\in\Delta_2$.

Fix any $0<\varepsilon<1$ and choose $0<d<a<1$. Since
$\varepsilon a<1$,
\[
\frac{M_{p,q}(\varepsilon a)-M_{p,q}(\varepsilon d)}
     {M_{p,q}(a)-M_{p,q}(d)}
=\varepsilon^{q+1}.
\tag{7}
\]
On the other hand, as $b\to\infty$,
\[
\frac{M_{p,q}(\varepsilon b)-M_{p,q}(\varepsilon a)}
     {M_{p,q}(b)-M_{p,q}(a)}
\longrightarrow \varepsilon^{p+1}.
\tag{8}
\]
Because $0<p<q$ and $0<\varepsilon<1$, we have
$\varepsilon^{q+1}<\varepsilon^{p+1}$. Thus, for all sufficiently large $b$,
condition \eqref{eq:increment} holds. Lemma \ref{lem:block} applies and gives
$L_{M_{p,q}}\notin\Xclass$.

\section*{Concrete sanity check}

The script \texttt{code/verify\_concrete.py} checks one explicit instance:
$p=1$, $q=2$, $\varepsilon=1/2$, $d=1/4$, $a=1/2$, $b=20$. It prints
\[
0.125=
\frac{M(\varepsilon a)-M(\varepsilon d)}{M(a)-M(d)}
<
\frac{M(\varepsilon b)-M(\varepsilon a)}{M(b)-M(a)}
\approx 0.2494,
\]
then constructs admissible $B,C,T$ and verifies the prefix and global modular
inequalities used in the proof.

\section*{Limitations and novelty check}

This is not a complete characterization of $\Xclass$. It gives a new explicit
negative family inside Orlicz spaces, complementing the source paper's
sufficient condition and positive examples. The result suggests that the
monotonicity of the scaling profile is close to the correct Orlicz-side
dividing line, but this packet proves only the negative direction for a broad
piecewise-power family.

On 2026-07-03, web searches for exact and close phrases around
\emph{Majorization revisited}, $\Xclass$, Orlicz spaces, Nazarov--Podkorytov,
and the ratio $M'(\varepsilon x)/M'(x)$ found the source arXiv page but no
later answer or matching Orlicz counterexample family. Novelty confidence is
therefore moderate, pending expert literature review.

\section*{Human-review recommendation}

The main point to check is Lemma \ref{lem:block}: the functions must satisfy
the exact $NP$ crossing condition, the common tail must not enter the prefix
comparison, and the Luxemburg modular normalizations must imply the stated
norm inequalities. The specialization to $M_{p,q}$ is then elementary.

\begin{thebibliography}{9}
\bibitem{ALM}
S. V. Astashkin, K. V. Lykov, and M. Milman,
\emph{Majorization revisited: Comparison of norms in interpolation scales},
arXiv:2107.11854, 2021.
\end{thebibliography}

\end{document}
